\documentclass[11pt]{article}

\usepackage[margin=25mm]{geometry}

\usepackage{amsmath,amssymb,fontspec,unicode-math}

\IfFontExistsTF{DejaVu Serif}{\setmainfont{DejaVu Serif}}{\setmainfont{Latin Modern Roman}}

\IfFontExistsTF{Latin Modern Math}{\setmathfont{Latin Modern Math}}{\setmathfont{DejaVu Math TeX Gyre}}

\usepackage{xurl,hyperref}

\setlength{\parindent}{0pt}\setlength{\parskip}{6pt}

\begin{document}

\title{Accelerated Evolving Set Processes for Local PageRank Computation}

\author{ELVIA preliminary AI analysis}\date{21 October 2026}\maketitle

Preliminary assessment: Minor.\par

Independent scientific review remains pending. Author adoption has not been requested.

The public local-PageRank routines have queue/reactivation defects, and their saved error metric and timing boundaries differ from the named comparisons. The abstract accelerated framework is not refuted.

\section*{Review scope}

Latest-version, claim-focused review of the cited author-paper passages and applicable current author-release paths. The stated counterexamples and scope controls were checked analytically; no author program, old verifier, training run or benchmark was executed.

\section*{Problem 1: The inner queue can stop before processing active vertices}

Source: Theorem A.3, PDF pp.18–20; Algorithm 4, PDF p.28; ppr\_solver.py, \_locAPPR, lines 253–333

Assessment: supported; impact: minor

\subsection*{Source claim (paraphrase)}

`ppr\_solver.py`, `\_locAPPR`, allocates a circular queue of one size but advances indices with a different modulus.

\subsection*{Analysis}

The queue capacity is \(n+1\), but increments use modulo \(n\) and sentinel detection uses an ordinary index difference. On a two-vertex single-edge graph, take \(\alpha=1/4\), \(\eta=1/2\), unit seed \(e_1\), center/start \(y=z_0=(2,2)^\top\), and tolerance \(0.01\). The inner quadratic has \(H=\begin{pmatrix}9/8&-3/8\\-3/8&9/8\end{pmatrix}\), \(b_{\rm eff}=(5/4,1)^\top\), and gradient \((1/4,1/2)^\top\). Both vertices are active. Queue writes place vertex zero in slot one, vertex one in slot zero, then overwrite slot one with the sentinel; the routine returns without updates. Exact minimizer is \(z^*=(19/12,17/12)^\top\), so \(h(z_0)-h(z^*)=\tfrac12(z_0-z^*)^\top H(z_0-z^*)=19/96>0.01\). This center satisfies the inner contract but is not asserted reachable in the default outer run.

\subsection*{Finding and scope}

The routine violates its inner accuracy contract on valid inputs. A default end-to-end convergence failure is not established.

\subsection*{Proposed correction}

Use the allocated capacity for every circular increment and circular distance test, or an explicit deque/sentinel design with consistent emptiness semantics.

An explicit local correction is described; author adoption and a corresponding published erratum are not established.

\section*{Problem 2: The lazy baseline drops a still-active processed coordinate}

Source: Section 4 baseline; Lemma 2.1, PDF p.3; ppr\_solver.py, appr, lines 26–99

Assessment: supported; impact: minor

\subsection*{Source claim (paraphrase)}

The baseline `appr` routine checks neighbors for activation after a push but can leave its own retained residual unqueued.

\subsection*{Analysis}

On a four-vertex star with a leaf source, \(\alpha=1/4\) and \(\epsilon=1/5\), the processed source retains residual \(3/8\). Its neighbor threshold prevents reactivation through the center. Returned estimate \((0,1/4,0,0)\) differs from exact PageRank \((3/8,19/40,3/40,3/40)\), with normalized error \(9/40>1/5\).

\subsection*{Finding and scope}

This baseline path fails its requested normalized tolerance. Other PageRank routines are not covered.

\subsection*{Proposed correction}

Requeue the processed coordinate whenever its retained residual remains active, preserving the same activation invariant as neighbors.

An explicit local correction is described; author adoption and a corresponding published erratum are not established.

\section*{Problem 3: The saved error is not the figure’s degree-normalized error}

Source: Figures 2 and 4 and convergence figures; ppr\_solver.py, norm\_inf and saved errs calls

Assessment: supported; impact: minor

\subsection*{Source claim (paraphrase)}

The evaluation labels a density-normalized infinity error, but its saved `norm\_inf` path does not divide the recovered PageRank error by degree.

\subsection*{Analysis}

For a center-seeded four-star and zero estimate, the saved error is \(5/8\), while the named degree-normalized error is \(5/24\). Nonconstant degrees make these different metrics rather than a universal common rescaling.

\subsection*{Finding and scope}

The public metric label and computation differ. A specific published ranking reversal is not established.

\subsection*{Proposed correction}

Use the declared degree normalization, or relabel and interpret the actual metric consistently.

An explicit local correction is described; author adoption and a corresponding published erratum are not established.

\section*{Problem 4: Reported timers cover unequal portions of the methods}

Source: Section 4, timing figures and Table 4; ppr\_solver.py, \_locAPPR, \_locGD and outer timer accumulation

Assessment: supported; impact: minor

\subsection*{Source claim (paraphrase)}

The accelerated path sums inner timers that start after initialization and scans, omitting outer residual/momentum work; the baseline timer begins before its initialization.

\subsection*{Analysis}

Those boundaries omit different necessary work, so the returned times cannot directly support an end-to-end speed ratio. This is a source-level accounting observation rather than a new latency measurement.

\subsection*{Finding and scope}

The public timing outputs need a common scope. No corrected large-graph speedup is estimated.

\subsection*{Proposed correction}

Time the same complete-call boundary for both methods, state warmup/synchronization rules and preserve component timings separately.

An explicit local correction is described; author adoption and a corresponding published erratum are not established.

\section*{Limits}

Preliminary AI assessment; independent human scientific adjudication remains pending.\par

An analytic counterexample or a release discrepancy does not by itself remeasure the paper’s reported experiments.\par

\section*{Source identity}

NeurIPS 2025 official proceedings PDF; compared with current arXiv:2510.08010v4

\url{https://proceedings.neurips.cc/paper_files/paper/2025/file/946ecab300b0695fe24b53a92e632935-Paper-Conference.pdf}

Reviewed PDF SHA-256: \nolinkurl{42d3b86c164a79cb7df29d93ceac61e6628d794555a1188f87c5de9bd20b1829}

The selected author/publisher source is the same work as the original case, matched by exact title, author identity and official version channel. The exact selected PDF SHA-256 was recomputed in the latest-version pass for Issue \#27. The PDF is a retained author/publisher copy; its original retrieval time is stated only when supported by matching acquisition metadata. File modification time was not substituted.

\end{document}
